\chapter{Quadratische Restwerte}
\label{b10guide:ch:quadratic-remainders}

Dieses Kapitel behandelt natürliche Lösungen quadratischer
Restgleichungen. Alle Zahlen und Lösungen liegen in $\mathbb N$;
die Aussagen verwenden ausschließlich natürliche Restwerte.

\section{Quadratische Restwerte und Eulers Kriterium}

\begin{definition}[Quadratischer Restwert]
\label{b10guide:def:quadratic-remainder}
Für $m\ge2$ und $a\in\mathbb N$ bezeichnet $Q_m(a)$ die Aussage
\[
 Q_m(a)\quad\Longleftrightarrow\quad
 \exists x\in\mathbb N\;
 \bigl(x<m\ \land\ x^2\mathbin{\mathrm{mod}}m
                       =a\mathbin{\mathrm{mod}}m\bigr).
\]
Ist $Q_m(a)$ erfüllt, heißt $a\mathbin{\mathrm{mod}}m$ ein
\emph{quadratischer Restwert zum Divisor $m$}.
Der Restwert $0$ ist ausdrücklich eingeschlossen.
\end{definition}

\begin{theorem}[Euler-Kriterium und Anzahl quadratischer Lösungen]
\label{b10guide:thm:euler-quadratic-criterion}
Seien $p$ eine ungerade Primzahl, $a\in\mathbb N$ und
$h=(p-1)\mathbin{\mathrm{div}}2$.
Für $p\nmid a$ gilt
\[
 Q_p(a)\quad\Longleftrightarrow\quad
 a^h\mathbin{\mathrm{mod}}p=1.
\]
Ist $p\nmid a$ und $Q_p(a)$ nicht erfüllt, so gilt
$a^h\mathbin{\mathrm{mod}}p=p-1$.
Die natürlichen Lösungen der Restgleichung
\[
 x^2\mathbin{\mathrm{mod}}p=a\mathbin{\mathrm{mod}}p,
 \qquad x<p,
\]
sind wie folgt bestimmt:
\begin{itemize}
 \item Für $p\mid a$ ist $x=0$ die einzige Lösung;
       außerdem gilt $a^h\mathbin{\mathrm{mod}}p=0$.
 \item Für $p\nmid a$ und erfülltes $Q_p(a)$ gibt es genau zwei
       Lösungen. Ist $r$ eine davon, so sind sie $r$ und $p-r$;
       dabei gilt $0<r<p$, und beide Lösungen sind verschieden.
 \item Für nicht erfülltes $Q_p(a)$ gibt es keine Lösung.
\end{itemize}
Vergleichsquelle: \cite{natquadraticbasic}.
\end{theorem}

\section{Ergänzungssätze und quadratische Reziprozität}

\begin{theorem}[Ergänzungssätze für quadratische Restwerte]
\label{b10guide:thm:quadratic-supplementary-laws}
Für jede ungerade Primzahl $p$ gilt
\[
 \begin{aligned}
 Q_p(p-1)&\quad\Longleftrightarrow\quad
               p\mathbin{\mathrm{mod}}4=1,\\
 Q_p(2)&\quad\Longleftrightarrow\quad
               p\mathbin{\mathrm{mod}}8=1
               \ \lor\ p\mathbin{\mathrm{mod}}8=7.
 \end{aligned}
\]
Vergleichsquellen: \cite{natquadraticbasic,natquadraticreciprocity}.
\end{theorem}

\begin{theorem}[Quadratisches Reziprozitätsgesetz]
\label{b10guide:thm:quadratic-reciprocity}
Seien $p$ und $q$ verschiedene ungerade Primzahlen.
Ist $p\mathbin{\mathrm{mod}}4=1$ oder
$q\mathbin{\mathrm{mod}}4=1$, so gilt
\[
 Q_p(q)\quad\Longleftrightarrow\quad Q_q(p).
\]
Sind hingegen $p\mathbin{\mathrm{mod}}4=3$ und
$q\mathbin{\mathrm{mod}}4=3$, so gilt
\[
 Q_p(q)\quad\Longleftrightarrow\quad \neg Q_q(p).
\]
Im zweiten Fall ist somit genau eine der beiden quadratischen
Restgleichungen lösbar \cite{natquadraticreciprocity}.
\end{theorem}
